\subsection{Broadband Demand and Price Elasticity}

The telecommunications demand literature has extensively studied price elasticity for internet access services, including early broadband adoption \citep{hausman2001private, rappoport2003demand}. Elasticity estimates vary substantially across methodologies and contexts. Market structure is a primary driver of this heterogeneity: \citet{cardona2009demand} find DSL own-price elasticity ranging from $-0.97$ in areas without platform competition to $-2.55$ where cable and mobile alternatives are available. Our EaP estimate falls at the lower end of the broader literature, consistent with EaP markets' intermediate level of platform competition and tighter income constraints, extending evidence to 2010--2024.

A key theoretical insight is that demand elasticity varies with necessity status. As connectivity integrates into economic and social life, price sensitivity declines as services shift from discretionary to essential \citep{becker1988theory}. Our finding of declining elasticity from 2015--2024 aligns with this theoretical framework.

Recent work emphasizes time-varying elasticity in technology markets. Evidence from broadband markets suggests that price sensitivity declines as services become embedded in daily economic life and as complementary technologies (smartphones, streaming, cloud applications) raise the opportunity cost of disconnection. Our contribution extends this literature by documenting elasticity evolution across a 15-year panel with explicit pre-trend testing.

Several studies published in this journal have examined the \textit{intensive margin} of broadband demand---consumers' willingness to pay for higher speeds or better quality conditional on being connected \citep{liu2018distinguishing, lindlacher2021low, sinclair2023broadband}. Our paper complements this strand by focusing on the \textit{extensive margin}: the adoption decision itself. Whether a household subscribes at all is governed by price relative to income and infrastructure availability---both of which we measure directly. The regional comparison of EU versus EaP countries provides heterogeneity that single-country intensive-margin studies cannot capture: income constraints bind far more tightly in EaP markets, so extensive-margin elasticity is substantially higher, while both regions converge to near-zero elasticity as broadband transitions from discretionary to essential service.

\subsection{Regional Heterogeneity and Development Context}

\citet{roller2001telecommunications} identify significant income effects on telecommunications adoption in cross-country panels, consistent with affordability constraints governing uptake where incomes are lower.

Our analysis focuses on Eastern Partnership countries---Armenia, Azerbaijan, Belarus, 
Georgia, Moldova, and Ukraine---which represent 
middle-income markets with substantial development 
heterogeneity relative to the EU \citep{ecorys2013evaluation}. Critically, 
this grouping is not a researcher's 
sampling choice but an institutionally and politically defined set, 
established under the EU's Eastern Partnership policy framework and 
subsequently operationalised through the EU4Digital initiative, 
which targets precisely these six countries for telecommunications regulatory 
harmonisation and progressive integration into the European Digital Single Market \citep{european2021eap}. The research question examined in this paper---whether demand-side convergence to EU necessity-good status is feasible over a policy-relevant horizon---is meaningful only for this institutionally defined group, as it concerns a specific integration trajectory rather than a random draw from a broader population. Consequently, $N{=}6$ reflects the exhaustive membership of the Eastern Partnership, not a statistical sample: there is no broader population of EaP countries from which additional observations could be drawn, and adding non-EaP transition economies would conflate distinct policy regimes. Limited prior work examines broadband demand elasticity in this region, making the present analysis a direct contribution to an understudied policy context.

\subsection{COVID-19 and Digital Transformation}

The pandemic's impact on telecommunications demand has received substantial attention, though primarily descriptive \citep{oecd2021covid, iab2020internet}. \citet{feldmann2021year} document more than 20\% traffic increases within weeks of lockdown announcements in March 2020, while \citet{favale2020campus} show usage patterns shifted dramatically toward videoconferencing and streaming. However, causal evidence on price elasticity changes remains limited.

Theoretically, COVID-19 could reduce elasticity through two mechanisms. 
First, lockdowns and social distancing converted broadband from convenience to 
necessity for remote work and education \citep{brynjolfsson2020covid}. 
Second, once connectivity became non-negotiable for school and employment, households continued subscribing even as incomes fell---treating broadband the way they treat electricity, so that the usual price-cutting response to financial stress no longer applied \citep{oecd2021covid}. Our contribution is econometric identification of elasticity changes using panel methods rather than descriptive evidence.

Importantly, the evidence presented here suggests digital transformation predated COVID-19. Our placebo tests show elasticity had already begun declining before 2020, consistent with survey evidence on the rapid pre-pandemic shift toward remote work \citep{brynjolfsson2020covid}, pointing to gradual evolution rather than a pandemic-induced exogenous shock.

Broadband policy encompasses supply-side interventions (infrastructure subsidies, spectrum allocation) and demand-side measures (affordability programmes, digital literacy) \citep{oecd2020broadband}. The effectiveness of demand-side policies depends critically on price elasticity: evidence on demand-side adoption programmes suggests mixed effectiveness \citep{hauge2010demand}, and where elasticity is near zero the case for price-based interventions weakens considerably \citep{guermazi2021digital}. Our evidence of declining elasticity supports the view that policy should shift from affordability to availability and quality as markets mature \citep{oecd2020broadband, ecorys2013evaluation}.
